\section{Discussion}
\label{sec:discussion}

The clearest result is not that spatial input always improves a VAE. It is that
the way a learned representation is explained matters strongly for peak
selection. Both centre-only and contextual GMM-targeted IG exceeded legacy
msiPL across the complete GBM and CAC evaluations, while first-layer L2 often
failed. This is consistent with the methodological difference: L2 observes one
linear transformation, whereas IG follows the full nonlinear path to a
label-free latent-space decision.

RQ1 has a conditional answer. Uniform context did not consistently improve GBM
reconstruction or mSCF1, but improved CAC mSCF1 in every paired section. A
neighbour mean can reduce local noise when adjacent spectra share relevant
structure, but can also blur boundaries or combine heterogeneous tissue. The
different spectral resolution, tissue classes, coverage and acquisition
properties of GBM and CAC plausibly change that balance. The experiments show
the boundary of the effect; they do not identify one biological cause.

RQ2 is supported more broadly. Nonlinear IG improved matched-count peak quality
over legacy msiPL across both collections. Completeness and deletion tests add
important evidence beyond mSCF1: the former checks numerical consistency with
the explained output, while the latter asks whether removing highly ranked bins
actually changes that output. Nevertheless, a high mask-correlation score does
not validate a peak as a biomarker. Independent molecular identification and
biological validation remain outside this study.

RQ3 shows a trade-off. Corrected attention gave a small but seed-consistent
mSCF1 improvement on the development section, yet did not improve deletion
faithfulness and consumed more memory. It should therefore remain an ablation
until tested across independent sections. Uniform averaging remains the more
defensible broad-validation model because it is parameter-free and was frozen
before whole-dataset evaluation.

S3PL is the strongest executable CAC baseline in peak quality. The present
method should not be described as both more accurate and cheaper than S3PL: its
matched CAC mSCF1 is lower, and the 10- versus 100-epoch protocols make the
recorded runtimes non-equivalent. The value of Spatial-msiPL is instead the
controlled evidence about context and attribution within a VAE framework.

Limitations include seed-1 section-wide validation, three-seed testing on only
one development section, sections that are not independent patients,
dataset-specific peak budgets, a GBM S3PL reproduction gap, and variable
cluster hardware availability. Future work should repeat training across
patients and sections, constrain attention more efficiently, investigate
automatic peak-budget selection, and validate selected ions biologically.
